\chapter{Motivation}
\label{chap:motivation}
\noindent \rule{6.6in}{0.01in}
\newpage

\section{Why Event Severity Matters}

It is worth being explicit about why this project treats severity estimation as a safety problem rather than only an efficiency one. Recurring congestion mostly costs commuters time -- an unpleasant but bounded loss. Event-driven congestion is different: an accident that is under-resourced can block the path an ambulance would otherwise take; a large gathering that is under-policed can turn a traffic problem into a crowd-safety one. The failure mode runs in both directions -- sending too few officers to a genuinely serious event is dangerous, but sending too many to a minor one quietly starves the rest of the city of manpower it may need at the same moment. A severity estimate that is both fast and reasonably accurate is therefore not a convenience; it is what makes the rest of the resource-allocation decision defensible.

\section{Limitations of Existing Approaches}

Searching the traffic-safety and traffic-engineering literature (surveyed properly in Chapter~\ref{chap:lit-survey}) makes the shape of the gap fairly clear. Most published machine-learning work on this general area falls into one of two camps: forecasting recurring flow or congestion levels with time-series or spatio-temporal deep models, or scoring the severity of a road-accident record after the fact, typically for post-hoc safety analysis or insurance-style reporting. Neither of these is quite the problem a traffic-control room actually faces in the moment an event is first logged, which is: given only what is known right now -- location, cause, time, whether the road is a major corridor -- how severe does this look, and what should we actually do about it. Very little of the published work couples a real-time, event-driven severity estimate directly to an operational recommendation of that kind, which is the gap this project tries to sit in: the feature set used here is deliberately restricted to fields an operator would have at the moment of logging, and the output is a concrete manpower/barricading/diversion plan rather than a bare probability.

\section{Why an Ensemble of Heterogeneous Models}

The event log used in this project is not a clean, single-type dataset -- it mixes continuous geographic coordinates, categorical descriptors of cause and corridor, and temporal flags derived from a timestamp, and there is no strong prior reason to expect one family of model to dominate on all of it. Gradient-boosted trees such as LightGBM and XGBoost tend to do well on this kind of structured, mixed-type tabular data because they split directly on thresholds and categorical levels without needing the inputs to be on a common scale. Neural approaches -- a plain multi-layer perceptron, or the attention-based TabNet -- can in principle pick up interaction effects between features (for instance, a particular event cause combined with a particular time bucket) that a sequence of axis-aligned tree splits approximates only indirectly. Rather than commit to one family in advance, this project trains one of each and combines their outputs with a stacked meta-learner, so the combination step -- not a hand-picked choice up front -- decides how much to trust each model.

\section{Why Class Imbalance Deserves Explicit Attention}

The single fact that shapes the evaluation methodology in this thesis more than any other is that severe events are rare. This is, on balance, a good thing about the world -- most reported traffic events genuinely are minor -- but it is an awkward thing for a classifier, because a model can learn to mostly ignore the rare Critical class and still post a very high overall accuracy simply by getting the common classes right. That failure mode is not hypothetical here; it shows up directly in the results in Chapter~\ref{chap:results}, where the model with the highest raw accuracy on this dataset turns out to be the worst at catching Critical events. Anticipating this before running a single experiment is why the evaluation in this thesis is built around macro-averaged F1 and per-class recall from the start, rather than being added afterwards to explain away a disappointing accuracy number.
